\documentclass[11pt]{article}
\usepackage[margin=0.82in]{geometry}
\usepackage{amsmath,amssymb,microtype}
\usepackage[hidelinks]{hyperref}

\title{The totally-positive Gabor frame-set conjecture is solved}
\author{Literature status for arXiv:1104.4894}
\date{13 August 2026}

\newcommand{\R}{\mathbb R}
\newcommand{\Z}{\mathbb Z}

\begin{document}
\maketitle

\begin{center}
\fbox{\parbox{0.91\textwidth}{\textbf{Status: literature already answered.}
Theorem 1 of arXiv:2608.04992 proves the precise frame-set conjecture from
arXiv:1104.4894 for every continuous integrable totally positive function
and for all real lattice parameters.}}
\end{center}

\section*{Source conjecture}

K. Gr\"ochenig and J. St\"ockler, \emph{Gabor Frames and Totally Positive
Functions}, arXiv:1104.4894, first prove the finite-type case.  In the
introduction (source lines 204--206) and again in ``Remarks and Conjectures''
(lines 1133--1140), they formulate the residual conjecture: if
$g\in L^1(\R)$ is continuous and totally positive, then its Gabor system
\[
 \mathcal G(g,\alpha,\beta)
 =\{e^{2\pi i\beta l t}g(t-\alpha k):k,l\in\Z\}
\]
is a frame for $L^2(\R)$ exactly when $\alpha\beta<1$.

The paper's two literal ``open problem'' extraction signals instead discuss
sharp nonuniform sampling and immediately announce the authors' own
finite-type answer.  The conjecture above is the distinct, precise residual
problem worth checking.

\section*{Exact later answer}

J. de Dios Pont, K. Gr\"ochenig, L. Liehr, I. Shafkulovska, and
M. A. Taylor, \emph{Gabor Frames of Totally Positive Functions: A Complete
Characterization}, arXiv:2608.04992 (submitted 5 August 2026), explicitly
identifies the earlier frame-set conjecture and says that it is proved in
full generality.  Its Theorem 1 states
\[
 \mathcal F(g)=\{(\alpha,\beta)\in\R_+^2:\alpha\beta<1\}
\]
for every continuous totally positive $g\in L^1(\R)$, equivalently
$\mathcal G(g,\alpha,\beta)$ is a frame if and only if
$\alpha\beta<1$.  This is an exact hypothesis-and-conclusion match.

\section*{Proof architecture checked in the resolving paper}

The nontrivial sufficiency direction is reduced by dilation to
$\mathcal G(g,\alpha,1)$ with $0<\alpha<1$.  The paper then:
\begin{enumerate}
\item extracts from every coset $x+\alpha\Z$ a sequence
  $\{k+\delta_k(x)\}$ whose perturbations lie in an interval avoiding the
  unique zero of the Zak transform and are recurrent under integer shifts;
\item uses total positivity and the de Boor--Friedland--Pinkus criterion to
  make the associated pre-Gramian
  $(g(k+\delta_k-l))_{k,l\in\Z}$ surjective on $\ell^1$;
\item upgrades surjectivity to Fredholmness, while recurrence makes the
  pre-Gramian a self-similar limit operator;
\item applies the limit-operator injectivity theorem, spectral invariance,
  and the Gabor submatrix criterion to obtain the frame property.
\end{enumerate}
The standard density obstruction gives necessity.  The paper also supplies
a Lean formalization of the new sufficiency direction.  Thus this packet
records literature priority and makes no novelty claim.

\end{document}
